\documentclass[11pt,a4paper]{article}
\usepackage[T1]{fontenc}
\usepackage{lmodern}
\usepackage[margin=25mm,headheight=14pt]{geometry}
\usepackage{amsmath,amssymb,mathtools,bm}
\usepackage{booktabs,longtable,tabularx,array}
\usepackage{enumitem,microtype,xcolor,fancyhdr,listings}
\usepackage[obeyspaces]{url}
\usepackage{hyperref}
\definecolor{inkblue}{HTML}{244B68}
\hypersetup{colorlinks=true,linkcolor=inkblue,urlcolor=inkblue,
 pdftitle={BundleFuse: the pipeline implemented in code},pdfauthor={BundleFuse implementation notes}}
\pagestyle{fancy}\fancyhf{}
\fancyhead[L]{\small BundleFuse: code-derived pipeline}
\fancyhead[R]{\small 3 November 2026}
\fancyfoot[C]{\thepage}
\setlength{\parindent}{0pt}\setlength{\parskip}{6pt}
\setlength{\emergencystretch}{2em}
\setlist{itemsep=3pt,topsep=4pt}
\newcolumntype{Y}{>{\raggedright\arraybackslash}X}
\def\UrlFont{\small\ttfamily}
\newcommand{\code}[1]{\nolinkurl{#1}}
\expandafter\def\expandafter\UrlBreaks\expandafter{\UrlBreaks\do\,\do\;}
\newcommand{\src}[1]{\par\begingroup\raggedright\textit{Code: }\code{#1}\par\endgroup}
\newcommand{\R}{\mathbb R}
\newcommand{\I}{\mathcal I}
\newcommand{\diag}{\operatorname{diag}}
\newcommand{\off}{\operatorname{off}}
\newcommand{\pos}[1]{\left[#1\right]_+}
\newcommand{\normrows}{\mathcal N}
\newcommand{\mean}{\operatorname{mean}}
\newcommand{\std}{\operatorname{std}}
\newcommand{\nnz}{\operatorname{nnz}}
\newcommand{\sig}[1]{\mathrm{#1}}
\lstset{basicstyle=\ttfamily\small,breaklines=true,columns=fullflexible,
 frame=single,framerule=.3pt,rulecolor=\color{inkblue},showstringspaces=false}
\numberwithin{equation}{section}
\title{\textcolor{inkblue}{BundleFuse}\\[5pt]
 The pipeline implemented in code\\[7pt]
 \large A step-by-step mathematical specification}
\author{Derived directly from the current working source}
\date{3 November 2026}
\begin{document}
\maketitle
\begin{abstract}
This document follows the executable pipeline: resolve configuration, construct data folds,
build fixed artifacts, compute and standardize candidate signals, select fusion settings on
validation queries, save a fitted checkpoint, predict, and evaluate. It describes both the
POG/Spotify path and the distinct retail/commerce path, as well as the optional extended
profile. Equations specify what the functions compute; interpretation is separated from
implementation facts. Earlier documentation, paper descriptions, and historical result tables
are not sources for this account. No accuracy or novelty claim is inferred from the formulas.
\end{abstract}

\textbf{Snapshot.} The branch is \code{reorganized}. Its base commit is
\begin{center}\code{162dd8ce8846b91bca3b508ede45fc6ebeaab576}\end{center}
The working tree contains changes,
so that commit alone does not identify this implementation. The companion
\code{pipeline_theory_sources.json} records SHA-256 hashes for the 30 inspected runtime,
configuration, entry-point, and timing-script files. Their combined source-set digest is
\begin{center}\small\ttfamily d9fe136d0d7a084d3c2e99510e5ea4acde35785b4515bf3a5e6018b8366b8033\end{center}
The source map in Appendix~\ref{app:map} connects the derivations to functions.

\clearpage
\tableofcontents
\clearpage

\section*{Notation and the execution contract}
\addcontentsline{toc}{section}{Notation and the execution contract}\label{sec:notation}
All core symbols are defined here before use. Additional symbols are introduced locally,
immediately before their equations. Bold lower-case letters are column vectors; matrices
are upper-case. A matrix row stores an item's features, and a score-matrix column stores
one query's scores over the entire catalogue.

\begin{longtable}{@{}>{\raggedright\arraybackslash}p{.21\linewidth}>{\raggedright\arraybackslash}p{.73\linewidth}@{}}
\toprule Symbol & Definition \\\midrule\endhead
$n$, $\I$ & Number of catalogue items and the internal ID set $\I=\{0,\ldots,n-1\}$. \\
$m$, $b_e$ & Number of training bundles and training bundle $e$, $1\le e\le m$. Different training rows may contain the same item set. \\
$S_j$, $G_j$ & Original revealed seeds and held-out ground-truth items for query $j$; $S_j\cap G_j=\varnothing$. \\
$S_j^\star$ & Seeds actually used for scoring after the optional seed cap. All original $S_j$ are excluded from recommendations. \\
$v$, $q$ & Number of validation queries and number of queries in a scoring batch. \\
$\bm x_j$, $X$ & Indicator $x_{ij}=\mathbf 1[i\in S_j^\star]$ and batch matrix $X=[\bm x_1,\ldots,\bm x_q]\in\R^{n\times q}$. \\
$B$, $B_0$ & Training bundle--item incidence and its binary incidence, both $m\times n$. In ordinary \code{BundleFuse.fit}, both are binary and equal. \\
$f_i$, $d_e$ & Training exposure of item $i$ and size of bundle $e$. \\
$\ell_i$, $a_i$ & Item category/artist label and, when available, album label. \\
$F$ & A generic dense item-feature matrix; $F_i$ is row $i$. A subscript identifies a particular feature family. \\
$C$, $M$, $U$ & Zero-diagonal item co-occurrence, its positive-PMI matrix, and row-normalized collaborative SVD features. \\
$r^t_{ij}$, $z^t_{ij}$ & Raw and catalogue-standardized score of signal key $t$ for item $i$, query $j$. \\
$\mathcal K$, $\mathcal P$ & Ordered selectable signal menu and selected prior-key list. \\
$w_t$ & Nonnegative, validation-selected weight of signal key $t$; zero keys are removed after pruning. \\
$\alpha_d,\gamma_v,\beta$ & Item-exposure exponent, training-bundle volume exponent, and incidence-surplus coefficient. Defaults: $0.5,0.5,1.5$. \\
$\alpha_b$ & Dense content-bridge mixing coefficient; default $0.5$. \\
$\alpha_h$, $\gamma_R$ & Validation-selected incidence-mixture and concordance coefficients. These differ from $\alpha_d,\alpha_b$ and the concordance kernel parameter. \\
$Y_{ij}$, $P_{ij}$ & Fused observed score and unscaled prior score. \\
$s_0$, $\kappa$ & Validation-frozen prior scale and exposure shrinkage ratio. \\
$\widehat\theta_{ij}$ & Final fused score before original revealed items are excluded. \\
$K$ & Requested recommendation or metric cutoff. \\
\bottomrule
\end{longtable}

\textbf{Basic operations.} Define $[x]_+=\max(0,x)$, elementwise product $\odot$,
$\diag(\bm u)$ as a diagonal matrix, and $\off(A)$ as $A$ with diagonal set to zero.
For dense rows define
\begin{equation}\tag{N.1}\label{eq:norm}
 \normrows(F)_i=\frac{F_i}{\lVert F_i\rVert_2+10^{-8}}.
\end{equation}
Sparse text rows use scikit-learn's exact L2 normalization, which leaves zero rows zero;
it is not the epsilon-normalization in \eqref{eq:norm}. Most final signal arrays are float32.
The formulas below describe real-valued arithmetic; where operation order or precision
matters, the actual code order is stated explicitly.

\textbf{What fitting means.} There is no gradient-trained neural network in this path.
Fitting constructs empirical graph statistics, SVD features and sometimes k-means clusters,
then searches validation-selected weights and routing settings. These are data-fitted
artifacts and selected parameters; calling the algorithm ``parameter-free'' would obscure
that distinction. Test ground truth is not passed into \code{BundleFuse.fit}.

\textbf{End-to-end order.}
\begin{enumerate}
\item Resolve a dataset definition, preprocessing profile, split protocol, and model configuration.
\item Load training bundles, validation queries, test queries, catalogue features, and context metadata.
\item Build training-incidence operators and the appropriate content/context artifacts.
\item Compute and standardize validation signal columns over all catalogue items.
\item Select weights, higher-order mixture, prior scale, and shrinkage on validation only.
\item Retain fitted state; optionally evaluate and save it through the CLI.
\item For a new query, compute only the dependencies of the frozen fusion, then rank.
\item Evaluate with the same ranking rule and the checkpoint's reconstructed data protocol.
\end{enumerate}
The fit CLI evaluates before saving by default; saving is skipped with \code{--no-save},
and evaluation is skipped with \code{--no-eval}. The conceptual distinction between
fitting and inference is unchanged by that ordering.

\section{Step 1: resolve configuration and dispatch}\label{sec:dispatch}
\src{bundlefuse/datasets/__init__.py; bundlefuse/config.py; evaluator.py}
The registry maps eight dataset names into three loader families. Shared methods implement
\code{datasets.load(name)}, \code{datasets.get(name).config(n_items)}, and
\code{datasets.build(name)}. The last returns an \emph{unfitted} model and loaded dataset.
The standard fit call is
\begin{lstlisting}
model, ds = datasets.build("pog_dense")
model.fit(ds.train, ds.val_pairs, **ds.fit_kwargs())
\end{lstlisting}
\code{fit_kwargs()} forwards features, labels, optional context inputs, validation users,
preprocessing-specific data, and a copied provenance specification. Dataset loading by
itself does not fit or save a model.

\begin{table}[htbp]\centering\small
\begin{tabularx}{\linewidth}{@{}p{.25\linewidth}YYY@{}}
\toprule Datasets & Loader family & Default protocol & Ordinary query limits \\\midrule
POG, POG Dense, Spotify, Spotify Sparse & BundleMLLM & Supplied files & All validation/test queries \\
Dunnhumby & Retail & Random resplit & 800 validation, 2,000 test \\
Electronic, Food, Clothing & BundleRec commerce & Random resplit & 800 validation, 2,000 test \\
\bottomrule
\end{tabularx}
\caption{Registry and fit-CLI defaults. Explicit zero or explicit Python \code{None} limits mean full folds. \code{validate.py} passes \code{None} by default and therefore uses full folds.}
\end{table}

There are \emph{two independent profile settings}. The \code{data_profile} chooses
preprocessing; the model's \code{signal_profile} chooses the scorer/selection recipe.
Both support \code{research} and \code{extended}. All eight current dataset JSON files
set \code{signal_profile="research"}; default dataset loading also chooses
\code{data_profile="research"}. In contrast, constructing \code{BundleFuseConfig}
directly defaults its signal profile to \code{extended}. A caller that disables JSON
configuration does not automatically recover the dataset JSON's research profile.

Configuration precedence is: dataset JSON (or fallback when no JSON is used), then non-null
caller overrides, then explicitly supplied CLI overrides. The catalogue size comes from
the loaded dataset. Current JSON files set collaborative SVD dimension 128,
$\alpha_d=\gamma_v=0.5$, $\beta=1.5$, $\alpha_b=0.5$,
\code{compressed=true}, and \code{content_rank=null}. Thus the current default does
\emph{not} apply an additional content projection, despite compressed checkpoint storage.
Other dataclass defaults include four coordinate-search rounds, 32-query validation blocks,
400 soft clusters, and a zero seed cap (use all seeds).

After the common graph is built, the special basket path is taken exactly when
\code{signal_profile == 'research'} and the preprocessing payload's kind is
\code{retail} or \code{bundlerec}. Otherwise fitting uses the general path. A research
signal setting without the basket payload therefore does not activate basket scoring.
Research baskets require their full family set; family ablations use the extended path.
They also reject a non-null content projection or $\alpha_b\ne0.5$ rather than silently
ignoring those overrides.

\section{Step 2: construct the data and query contract}\label{sec:data}
\src{bundlefuse/datasets/base.py; bundlefuse/datasets/baskets.py; bundlefuse/splits.py}

\subsection{Supplied POG/Spotify folds}
\code{ProvidedSplits} reads training rows and paired input/ground-truth files for validation
and test. It checks catalogue size, item bounds, unique row IDs, nonempty rows, no duplicate
items inside a row, matching input/target query IDs, disjoint fold IDs, and no seed/target
overlap within a query. Query pairs are emitted in ascending original bundle-ID order.
Distinct training IDs with identical item contents are retained. Supplied mode neither
regenerates masks nor deduplicates bundles. Non-default resplit seeds or mask fractions
are rejected in supplied mode.

For non-music data, available title, image and description arrays are loaded separately.
For music, the available content tensor is loaded under the audio key. Each modality is
row-normalized by \eqref{eq:norm}; their concatenation is also supplied as clustering
input. Fashion labels come from \code{cate}/\code{cate_id}; music labels are artist
indices with separate album indices. Music is recognized by dataset name or artist metadata.

\subsection{Random resplit and reveal masking}
Let $M_b$ denote the number of retained complete bundles and let $r$ be the requested split
seed. Resplitting shuffles their index/order list with Python's \code{random.Random(r)};
training uses the first $\lfloor0.7M_b\rfloor$ bundles, validation the next
$\lfloor0.1M_b\rfloor$, and test the remainder. For POG/Spotify resplit mode, complete
bundles are first reconstructed from supplied inputs plus targets, deduplicated by sorted
item set, and restricted to at least two distinct items.

Let $\phi\in(0,1)$ be the hidden fraction, default $0.5$, and let $\widetilde b$ be a shuffled
copy of a held-out complete bundle. Before forming each query define the reveal count
\begin{equation}
 h(b)=\max\{1,\operatorname{round}((1-\phi)|b|)\}.
\end{equation}
The first $h(b)$ shuffled items form $S$, and the rest form $G$. If the latter is empty,
the last item is hidden and the others are revealed. Python rounding, including its
half-to-even rule, is used. A single RNG is reused across all bundles in each masking call;
validation uses seed $r+1$, test $r+2$. Positive query caps retain a prefix after masking;
zero/\code{None} retain the full fold. A query cap is not a new random sample.

\subsection{Default Dunnhumby preprocessing}
\begin{enumerate}
\item Group transaction rows by basket ID into unique product sets, retaining household and day metadata.
\item Keep baskets of 3--20 products; shuffle once with seed zero and cap at 30,000 baskets.
\item Count product appearances in that capped pool, retain products occurring at least twice,
      remove other products, and discard resulting baskets shorter than three items.
\item Map retained products, sorted by original ID, to contiguous internal IDs. Globally
      deduplicate identical mapped baskets, keeping the first household/day association.
\item Read product text from department, commodity, subcommodity and brand. Subcommodity
      gives a category; a product with missing metadata receives its own missing category.
\item Fit catalogue TF--IDF and SVD content features as described below, then resplit and mask.
      Household counts used by the model are built from the training fold only.
\end{enumerate}
The day column is read, but this loader uses a random split, not a chronological split.

\subsection{Default commerce preprocessing}
\begin{enumerate}
\item Read bundle--item rows, build the catalogue from all listed items, deduplicate items
      within bundles, retain bundles of at least two items, and globally deduplicate equal item sets.
\item Retain the first original bundle key and load its user and session associations.
\item Read the external user--item pretraining file and restrict its items to this catalogue.
      Its weighted co-occurrence graph is defined in Section~\ref{sec:contexts}.
\item Parse the first bracketed category path. The entire parsed path defines the leaf label;
      its first up-to-two entries define a second-level label. Missing labels are unique per item.
\item Construct text by concatenating the comma-joined title fields and parsed category path.
      Build TF--IDF/SVD features, then shuffle, resplit and mask the retained bundles.
\end{enumerate}
User-bundle timestamps are read but not used for a temporal restriction. User/session
co-occurrence from bundle interactions uses training bundles only; the external pretraining
file is loaded as supplied, with no timestamp filtering in this code.

\subsection{Text features and information available before splitting}\label{sec:text}
Let $c_{it}$ be the count of retained token/ngram $t$ in item document $i$, and let
$\operatorname{df}_t$ be the number of catalogue documents containing $t$. The configured
vectorizer uses English stop words, unigrams and bigrams, and sublinear term frequency.
With the installed library's smoothed-IDF defaults, the pre-normalized text weight is
\begin{equation}
 t_{it}=\begin{cases}
 (1+\log c_{it})\left(1+\log\dfrac{n+1}{\operatorname{df}_t+1}\right),&c_{it}>0,\\
 0,&c_{it}=0.
 \end{cases}
\end{equation}
Let $T\in\R^{n\times d_t}$ be the resulting sparse L2-normalized text matrix, where
$d_t$ is its retained vocabulary size. The vocabulary cap is 8,000 for retail and 12,000
for default commerce. Define $r_t=\min(128,d_t-1)$, and let $V_t$ denote the right
singular vectors returned by \code{TruncatedSVD(r_t, random_state=0)}. Dense basket
content is $F_c=\normrows(TV_t)$.

This preprocessing is catalogue-transductive: text vocabulary/SVD use all retained item
metadata before the bundle split. Retail's frequency filter also examines the pre-split
basket pool. This differs from fitting graph co-occurrence on training bundles. The code
therefore supports a known catalogue, not a claim that every preprocessing decision is
estimated from the training fold alone. Predictions use contiguous internal IDs; the
loaders do not export a dedicated original-ID mapping inside the fitted model.

\section{Step 3: build the common training graph}\label{sec:graph}
\src{bundlefuse/operators.py: SimplicialGraphOperator}
In ordinary model fitting define $B_{ei}=\mathbf1[i\in b_e]$ and $B_0=B$; repeated items
inside a training row are removed by the graph constructor. The helper also accepts bundle
weights, but \code{BundleFuse.fit} does not pass them. The equations in this section apply
to the binary incidence used by that public fit path. Define
\begin{align}
 f_i&=\sum_e B_{ei},&d_e&=\sum_i(B_0)_{ei},\\
 D&=\diag((f_i+5)^{-\alpha_d}),&
 V&=\diag((d_e+1)^{-\gamma_v}),&
 E&=\diag((d_e+1)^{-1}).
\end{align}
$f_i$ counts appearances in retained \emph{training rows}, so retained duplicate rows affect
exposure. The additive constants 5 and 1 are implementation choices, not symbolic limits.

\subsection{Co-occurrence, positive PMI and collaborative SVD}
Define $g_i=\sum_j C_{ij}+10^{-9}$ and $z_C=\sum_{i,j}C_{ij}$ after computing
\begin{equation}
 C=\off(B^\top B).
\end{equation}
For stored nonzero co-occurrences define
\begin{equation}\label{eq:itempmi}
 M_{ij}=\pos{\log\left(\frac{C_{ij}z_C}{g_i g_j}+10^{-12}\right)};
 \qquad M_{ij}=0\ \text{when }C_{ij}=0.
\end{equation}
The denominator uses co-occurrence row masses $g_i$, not bundle exposure $f_i$.
Let $r_c=\min(\text{\code{svd_dim}},\max(1,n-2))$ and $V_c$ be the truncated right singular
basis of $M$. The stored collaborative feature matrix is
\begin{equation}
 U=\normrows(MV_c).
\end{equation}
The SVD uses random seed zero; if $M$ has no nonzeros, $U$ is a zero matrix of the requested
clamped width. The feature is the transformed matrix (left singular vectors times singular
values), followed by row normalization, not the unscaled left singular vectors alone.

\subsection{The zero-diagonal propagation operator}
Define the symmetric propagation operator $L\in\R^{n\times n}$ by
\begin{equation}\label{eq:L}
 L=DB^\top EBD-\diag\left(D_{ii}^2\sum_eB_{ei}E_{ee}\right).
\end{equation}
Since $B$ is binary, the subtracted vector is the diagonal of the preceding product.
\code{propagate(A)} computes $LA$ by multiplying in the order
$B(DA)$, then $B^\top E[B(DA)]$, then left-scaling by $D$ and subtracting the diagonal
correction. It accepts a score vector or an $n$-row dense feature/score matrix and returns
float32. It does not use row-stochastic normalization.

\subsection{Two different pairwise objects}\label{sec:pairdistinction}
For a query indicator $\bm x$ define its training-bundle overlap vector
$\bm o=VB D\bm x$. The ordinary \code{pair} score is
\begin{equation}\label{eq:pair}
 \bm p=B^\top\bm o=B^\top VB D\bm x.
\end{equation}
Separately, the graph stores
\begin{equation}\label{eq:W}
 W=\off(B^\top VB D),
\end{equation}
under \code{W_pair}. $W$ is used for concordance; \eqref{eq:pair} is used for the
\code{pair} signal. They differ on the diagonal. Although revealed candidates are excluded
at the final ranking step, their scores participate in standardization and diffusion.
Thus replacing \eqref{eq:pair} by $W\bm x$ would change the implemented pipeline.

\subsection{Pure incidence surplus}
Define overlap counts $\bm k=B_0\bm x$ and let $\bm\delta$ denote the incidence-surplus
score vector. The code computes
\begin{equation}\label{eq:delta}
 \bm\delta=B^\top\bigl[\bm o\odot\beta\pos{\bm k-\bm1}\bigr].
\end{equation}
For candidate $i$, this adds evidence from training bundles that contain $i$ and at least
two revealed seeds. It is exactly zero for a single-seed query or whenever no training
bundle overlaps two scoring seeds. This is excess over a linear overlap score; the code
never requires all revealed seeds to occur in the same training bundle.

\section{Step 4: build dense content and ontology artifacts}\label{sec:features}
\src{bundlefuse/model.py: BundleFuse.fit; bundlefuse/operators.py: DualSpectrumOntology, BilateralContentBridge}
This section describes the general POG/Spotify path. Research basket content/context
artifacts are specified separately in Section~\ref{sec:contexts}.

\subsection{Input normalization, category labels and style residuals}
Let $F^{(1)},\ldots,F^{(b)}$ be the available content modalities, with $b$ denoting the
number of modalities. The loader normalizes each modality. A research full-ranking
preprocessing payload marks them as already normalized; otherwise the model normalizes
them again. Their raw content feature is
\begin{equation}
 F_{\mathrm{content}}=\normrows\bigl([F^{(1)}|\cdots|F^{(b)}]\bigr).
\end{equation}
If explicit category labels are absent, the model obtains them by k-means on supplied
clustering features, concatenated modalities, or collaborative features, in that priority.
The effective cluster count is $k_s=\max(2,\min(\text{\code{n_soft_clusters}},n-1))$.
Every k-means call uses \code{MiniBatchKMeans}, random seed zero, three initializations,
and batch size 4,096.

For each modality, define a style-group label $c_i$. For fashion, $c_i=\ell_i$; for music,
$c_i$ is a k-means cluster of audio/content features, not the artist label. For each group
$c$ define its catalogue centroid $\overline F_c$, then its normalized residual feature
\begin{equation}\label{eq:style}
 \overline F_c=\frac{1}{|\{i:c_i=c\}|}\sum_{i:c_i=c}F_i,
 \qquad F_{\mathrm{style},i}
 =\frac{F_i-\overline F_{c_i}}{\lVert F_i-\overline F_{c_i}\rVert_2+10^{-8}}.
\end{equation}
The keys are \code{style}, \code{imgstyle}, \code{descstyle}, and \code{audiostyle}
for title/content, image, description, and audio respectively. These are deterministic
centroid residuals after clustering/label construction, not outputs of a neural encoder.

\subsection{Dense bridge and optional content projection}
Let $F_s$ be the normalized concatenation of available style residuals. In the research
music path, $F_s$ is the audio style residual itself, with no extra normalization.
Define the propagated and bridged features
\begin{equation}\label{eq:bridge}
 F_p=\normrows(LF_s),\qquad
 F_{\mathrm{bridge}}=\normrows(F_s+\alpha_b F_p).
\end{equation}
This mixes item content with graph-neighbor content once. It is built during fitting and
reused at prediction time.

Only when both \code{compressed=true} and \code{content_rank=r} is non-null does the model
apply an additional rank-$r$ truncated-SVD transform to any content/style/bridge feature
whose input width exceeds $r$. Each eligible feature is projected separately, using seed
zero. Projected rows are \emph{not} normalized again. Collaborative $U$ is excluded from
this projection. The concordance descriptor is retained separately before projection,
so projecting the scoring feature table does not necessarily project concordance features.
With the current JSON defaults, this entire additional projection step is skipped.

\subsection{Category complementarity and saturation}
Let $c$ index the $n_c$ distinct category labels, and define the training category-count
matrix $A\in\R^{m\times n_c}$ by $A_{ec}=|\{i\in b_e:\ell_i=c\}|$.
Define category masses $a_c=\sum_e A_{ec}+10^{-9}$, total $a_+=\sum_c a_c$, and
$H=\off(A^\top A)$. For nonzero $H_{cd}$ define
\begin{equation}\label{eq:catpmi}
 Q_{cd}=\pos{\log\left(\frac{H_{cd}a_+}{a_c a_d}+10^{-12}\right)};
 \quad Q_{cd}=0\text{ otherwise}.
\end{equation}
The marginal in this category formula is a count of category appearances, not the row sum
of $H$. Different items with the same label contribute repeated counts to $A$.

For each category $c$, define $u_c$ as the number of training bundles containing that
category at least once and $v_c$ as the number containing it at least twice. Define the
smoothed duplication rate $\rho_c=(v_c+0.5)/(u_c+1)$. The non-music saturation score is
\begin{equation}\label{eq:saturation}
 r^{\sig{slot}}_i(S^\star)=
 -\sum_{s\in S^\star}\mathbf1[\ell_i=\ell_s]\,
 \mathbf1[\rho_{\ell_s}<0.15]\,(1-\rho_{\ell_s}).
\end{equation}
It penalizes a candidate sharing a rarely duplicated category with each matching seed.
The general fashion category signal adds one half of this score. General music and the
research basket category signals do not add this term.

\subsection{Soft ontology for the general non-music path}
Let $F_o$ be the concatenated modality input used by the ontology. The research fashion
path row-normalizes this concatenation; the extended path uses the concatenation directly.
Let $\bm\mu_c$ be k-means centre $c$, normalized to unit length with epsilon $10^{-8}$,
for $1\le c\le k_s$. Let $\tau_s$ be \code{soft_tau}, default $0.05$.
The soft membership matrix $J\in\R^{n\times k_s}$ is
\begin{align}
 t_{ic}&=F_{o,i}\bm\mu_c^\top/\tau_s,
 &J_{ic}&=\frac{\exp(t_{ic}-\max_d t_{id})}
 {\sum_d\exp(t_{id}-\max_d t_{id})+10^{-9}}.
\end{align}
Define each training bundle's soft sum $\bm b_e^s=\sum_{i\in b_e}J_i$ and the matrix
$B_s$ with those rows. The off-diagonal nonnegative soft co-occurrence is
\begin{equation}
 H_s=\left[\off\left(B_s^\top B_s-
             \sum_e\sum_{i\in b_e}J_i^\top J_i\right)\right]_+.
\end{equation}
The subtracted term removes each item's self-outer-product inside a training bundle.
Define catalogue soft masses $a^s_c=\sum_i J_{ic}$ and $a^s_+=\sum_c a^s_c$. The
stored soft compatibility matrix is computed over all entries as
\begin{equation}
 (Q_s)_{cd}=\pos{\log\left(
 \frac{(H_s)_{cd}a^s_+}{a^s_c a^s_d+10^{-9}}+10^{-12}\right)}.
\end{equation}
The model stores $J$ and $JQ_s$. The soft signal later matches a candidate's soft
compatibility row against the \emph{sum} of seed membership rows. This is separate from
the category-count PMI in \eqref{eq:catpmi}. If the ontology has no clustering features,
these soft artifacts are absent and \code{softcat} is unavailable.

\section{Step 5: build the research basket context artifacts}\label{sec:contexts}
\src{bundlefuse/datasets/baskets.py; bundlefuse/basket_signals.py: BasketSignals}
The research basket path returns before the general ontology/style construction. Its
content $F_c$ is the catalogue text-SVD matrix from Section~\ref{sec:text}. Its category
PMI uses the count construction of \eqref{eq:catpmi}, without slot saturation or soft labels.

\subsection{Retail household context and bridge}
Let $Z_h$ be the household-by-item training count matrix; $(Z_h)_{ui}$ counts appearances
of item $i$ in training baskets belonging to household $u$. Let
$d^h_u=\sum_i(Z_h)_{ui}+1$. Define
\begin{equation}
 \widetilde Z_h=\diag((d^h_u)^{-1/2})Z_h,
 \qquad C_h=\off(\widetilde Z_h^\top\widetilde Z_h).
\end{equation}
The retail household signal will use $C_hD\bm x$, including the common seed-exposure
weight $D$. Its dense bridge is
\begin{equation}
 F_h=\normrows\bigl(F_c+0.5\normrows(LF_c)\bigr).
\end{equation}
Household context is shared evidence across the training fold. It is not a per-query
household-personalization term in this research basket path; supplied query users do not
change \code{BasketSignals.score}.

\subsection{Commerce user/session graphs and external pretraining graph}
Let $Z_u$ and $Z_s$ be user-by-item and session-by-item count matrices built only from
training bundles and their recorded associations. Missing negative associations are skipped.
Define $C_u=\off(Z_u^\top Z_u)$ and $C_s=\off(Z_s^\top Z_s)$; these two contexts
are not degree-normalized here.

Let $Z_p$ be the external pretraining user-by-item count matrix after catalogue restriction,
and define $d^p_u=\max(\sum_i(Z_p)_{ui},1)$. Define its degree-normalized co-occurrence
and the commerce diffusion operator $L_c$ as
\begin{align}
 C_p&=\off\left(Z_p^\top\diag((d^p_u)^{-1})Z_p\right),\\
 D_c&=\diag((f_i+1)^{-1/2}),\\
 L_c&=D_c(C+0.5C_s+0.5C_u+0.2C_p)D_c.\label{eq:Lcommerce}
\end{align}
$L_c$ replaces $L$ for commerce score diffusion. It does not replace $L$ inside the
common graph constructor or the separately stored concordance matrix $W$.

Using sparse text $T$ and dense text-SVD content $F_c$, define
\begin{align}
 T_b&=\operatorname{L2rows}(T+1.5L_cT+0.5L_c(L_cT)),\\
 F_b&=\normrows(F_c+1.5L_cF_c+0.5L_c(L_cF_c)).
\end{align}
$T_b$ is the selectable sparse bridge; $F_b$ is a dense descriptor for concordance.
The two are different artifacts, and the dense one is not an additional selectable
\code{content} signal in the default commerce menu.

For commerce taxonomy, build three versions of \eqref{eq:catpmi}: leaf-category
$Q_\ell^{\mathrm{off}}$ with diagonal removed, leaf-category $Q_\ell^{\mathrm{keep}}$
with diagonal retained, and second-level $Q_2^{\mathrm{keep}}$ with diagonal retained.
In each version the appropriate label-count matrix supplies its own marginals. The
last two support the commerce homophily signal.

\section{Step 6: compute raw candidate signals}\label{sec:signals}
All formulas here use scoring seeds $S^\star$ and indicator $\bm x$. They produce one
score for \emph{every} catalogue item, including seeds; seed exclusion occurs only when
ranking. Batch implementations perform these same operations for columns of $X$.

\subsection{Shared collaborative, diffusion and content signals}
Define $\bm p$ and $\bm\delta$ by \eqref{eq:pair} and \eqref{eq:delta}. In commerce,
first replace the pair score by
\begin{equation}
 \bm p_c=\bm p+0.35C_u\bm x+0.35C_s\bm x.
\end{equation}
The family-specific raw score definitions are then:
\begin{center}\small
\begin{tabularx}{\linewidth}{@{}p{.20\linewidth}YY@{}}
\toprule Code key & General / research retail & Research commerce \\\midrule
\code{pair} & $\bm p$ & $\bm p_c$ \\
\code{simp} & $\bm p+\bm\delta$ & $\bm p_c+\bm\delta$ \\
\code{diff2_pair} & $L\bm p$ & $L_c\bm p_c$ \\
\code{diff2} & $L\bm p+L\bm\delta$ & $L_c(\bm p_c+\bm\delta)$ \\
\code{cooc} & $C\bm x$ & $C\bm x$ \\
\code{ppmi} & $M\bm x$ & $M\bm x$ \\
\code{pop} & $\log(1+f_i)$ for candidate $i$ & Not in the selectable menu \\
\code{csvd} & $U_i\bigl(|S^\star|^{-1}\sum_{s\in S^\star}U_s\bigr)^\top$ & Same \\
\bottomrule
\end{tabularx}
\end{center}
The code evaluates $L\bm p$ and $L\bm\delta$ separately before adding for general/retail
\code{diff2}. Algebraically it equals $L(\bm p+\bm\delta)$, but the prescribed float32
operation order is retained. Commerce uses the latter summed-input form with $L_c$.

For any selectable dense feature $F_t$, the raw content score is
\begin{equation}\label{eq:densescore}
 r^t_i(S^\star)=F_{t,i}\left(\frac{1}{|S^\star|}\sum_{s\in S^\star}F_{t,s}\right)^\top.
\end{equation}
This applies to raw content, styles, dense bridges, and collaborative SVD features when
they are in the selected menu. Research retail uses $F_c$ for \code{content} and $F_h$
for \code{bridge}. Research commerce instead scores
\begin{equation}
 \bm r^{\sig{tfidf}}=T(T^\top\bm x),\quad
 \bm r^{\sig{bridge}}=T_b(T_b^\top\bm x),\quad
 \bm r^{\sig{pretrain}}=C_p\bm x.
\end{equation}
Sparse text uses a \emph{sum} of seed columns; dense \eqref{eq:densescore} uses a mean.
Retail also computes TF--IDF in its full reference scorer, but that key is not selectable
in its default menu. Placeholder \code{contrast}/\code{img} values in the basket full
scorer are \code{None}, not trained neural components.

\subsection{Taxonomy, homophily and household scores}
For the relevant category matrix $Q$, the basic category score is
\begin{equation}
 r^{\sig{cat}}_i(S^\star)=\sum_{s\in S^\star}Q_{\ell_s,\ell_i}.
\end{equation}
General non-music scoring adds $0.5r_i^{\sig{slot}}$ from \eqref{eq:saturation}.
General music uses artist labels without that penalty. Commerce uses
$Q_\ell^{\mathrm{off}}$ without the penalty. Its homophily score, using second-level
label $h_i$, is
\begin{equation}
 r^{\sig{homo}}_i(S^\star)=\sum_{s\in S^\star}
 \left[(Q_\ell^{\mathrm{keep}})_{\ell_s,\ell_i}
       +0.5(Q_2^{\mathrm{keep}})_{h_s,h_i}\right].
\end{equation}
For soft ontology, the score is
\begin{equation}
 \bm r^{\sig{softcat}}=(JQ_s)\left(\sum_{s\in S^\star}J_s\right)^\top.
\end{equation}

For general music, let $H_a$ and $H_b$ be the item-by-artist and item-by-album binary
membership matrices. Let their column sizes be $n^a_c$ and $n^b_c$. Define
$\widetilde H_a=H_a\diag((n^a_c+5)^{-1/2})$ and
$\widetilde H_b=H_b\diag((n^b_c+2)^{-1/2})$. Then
\begin{equation}
 \bm r^{\sig{homo}}=\widetilde H_a\widetilde H_a^\top\bm x
                   +2\widetilde H_b\widetilde H_b^\top\bm x.
\end{equation}
The album term is zero if album labels are absent. Research retail's separate household
signal is $\bm r^{\sig{hh}}=C_hD\bm x$, code key \code{hh_cooc}.

\subsection{Concordance: support from pairs of revealed seeds}\label{sec:rao}
The concordance routine uses $W$ in \eqref{eq:W}, not the \code{pair} score matrix.
Let $F_R$ be the descriptor used only for this routine, $\eta$ the seed-kernel parameter,
$\lambda_d$ the dense-support coefficient, and $\tau_R=0.05$ its denominator stabilizer.
Before defining the kernel, define candidate-to-seed support and total support by
\begin{equation}
 q_{is}=W_{is}+\lambda_d[F_{R,i}F_{R,s}^\top]_+,
 \qquad A_i=\sum_{s\in S^\star}q_{is}.
\end{equation}
If no descriptor is available, omit the dense term. For seeds $s,t$, define similarity
$G_{st}=0.5\mathbf1[\ell_s=\ell_t]+0.5[F_{R,s}F_{R,t}^\top]_+$ when both labels and
features are available. With only one source use that source alone; with neither, the
routine starts with the identity. Its diagonal is set to one. Define the off-diagonal
seed kernel $K^S$ by
\begin{equation}
 K^S_{st}=\begin{cases}
 0,&s=t,\\
 [1+\eta G_{st}]_+,&s\ne t,\ \eta<0,\\
 1+\eta G_{st},&s\ne t,\ \eta\ge0.
 \end{cases}
\end{equation}
The raw concordance score is
\begin{equation}\label{eq:rao}
 r^{\sig{rao}}_i(S^\star)=
 \frac{\displaystyle\sum_{s\in S^\star}\sum_{t\in S^\star}
                 q_{is}K^S_{st}q_{it}}
      {A_i+\tau_R}.
\end{equation}
The double sum includes both ordered directions of each distinct seed pair; it is not
halved. It vanishes for fewer than two scoring seeds. For $A_i>0$, setting
$p_{is}=q_{is}/A_i$ rewrites \eqref{eq:rao} as
$A_i^2(A_i+\tau_R)^{-1}\sum_{s\ne t}K^S_{st}p_{is}p_{it}$.
This equivalent form explains agreement weighting; the implementation uses the quadratic
matrix expression directly, without dividing by $A_i$.

\begin{center}\small
\begin{tabularx}{\linewidth}{@{}p{.29\linewidth}Ycc@{}}
\toprule Path & Descriptor $F_R$ & $\eta$ & $\lambda_d$ \\\midrule
General fashion & Title/content style residual & $-1$ & $0.15$ \\
General music & Audio style residual & $+1$ & $0.15$ \\
Research retail & Dense retail bridge $F_h$ & $-1$ & $0.15$ \\
Research commerce & Dense commerce bridge $F_b$ & $-1$ & $0.25$ \\
\bottomrule
\end{tabularx}
\end{center}
The general full scorer creates this signal only when its descriptor exists. The selected
scorer requests it only when the frozen fusion needs it. $\eta$ inside the kernel is
\emph{not} the validation-selected multiplier $\gamma_R$ later applied to the entire
standardized concordance score.

\section{Step 7: standardize signals and fix their ordered menu}\label{sec:normalization}
\src{bundlefuse/operators.py: zcols; bundlefuse/model.py; bundlefuse/basket_signals.py}
For every signal $t$ and query $j$, define the catalogue mean
$\mu^t_j=n^{-1}\sum_i r^t_{ij}$. Define $\nu$ as the standard-deviation correction:
$\nu=1$ on the general Torch path and $\nu=0$ on the NumPy research basket path. Then
\begin{equation}\label{eq:z}
 z^t_{ij}=\frac{r^t_{ij}-\mu^t_j}
 {\sqrt{\frac{1}{n-\nu}\sum_{i'}(r^t_{i'j}-\mu^t_j)^2}+10^{-9}}.
\end{equation}
This normalization is over catalogue items \emph{within each query}, not over other queries,
not just warm items, and not after seeds have been masked. Constant columns become zero.
The Python types determine which standard deviation the shared \code{zcols} uses.

The auxiliary keys \code{pair}, \code{diff2_pair} and \code{rao} are standardized
separately. In particular, $z^{\sig{simp}}$ is the normalization of
$\bm p+\bm\delta$, not $z^{\sig{pair}}+z^{\sig{delta}}$. This matters for fusion.

The research menus below are ordered; coordinate selection visits keys in this order.
Unavailable feature keys are removed on the general path, and family filtering is applied.
Research basket menus are fixed and require their associated families.
\begin{description}[style=nextline]
\item[POG / POG Dense]
\code{cat, softcat, style, imgstyle, bridge, simp, diff2, pop, cooc, ppmi, csvd}.
Description style may contribute to the bridge, but is not a separate selectable key here.
Raw \code{content} is also excluded from this research menu.
\item[Spotify / Spotify Sparse]
\code{simp, diff2, pop, cooc, ppmi, csvd, cat, homo, audiostyle, bridge}.
\item[Research retail]
\code{simp, diff2, pop, cooc, hh_cooc, ppmi, csvd, cat, content, bridge}.
\item[Research commerce]
\code{simp, diff2, cooc, ppmi, csvd, cat, homo, tfidf, bridge, pretrain}.
\end{description}
The auxiliary keys are not themselves coordinate-selected weights. They are dependencies
of the higher-order mixture. On the extended path the menu instead follows the full
scorer's computed dictionary order, excluding auxiliary keys and applying the configured
family filter. Appendix~\ref{app:extended} specifies those additional signals.

\section{Step 8: select the frozen fusion on validation queries}\label{sec:selection}
\src{bundlefuse/fusion.py; bundlefuse/validation_cache.py; bundlefuse/operators.py: EmpiricalBayesRouter}
The selection routine needs a precise objective, score construction and search order.
This section defines them in that dependency order. Validation labels affect metric-based
choices; graph/content artifacts remain fixed during this search.

\subsection{Define the ranking rule and validation objective first}\label{sec:objectives}
Let $A^{\mathrm{score}}\in\R^{n\times v}$ denote any candidate validation score matrix.
For query $j$, let $\pi_{j1},\pi_{j2},\ldots$ be the items outside original $S_j$, ordered
by decreasing $A^{\mathrm{score}}_{ij}$ and then increasing internal item ID for exact ties.
Define $h_{jr}=\mathbf1[\pi_{jr}\in G_j]$ and
$K_j=\min(K,n-|S_j|)$. For nonempty ground truth, the per-query metrics are
\begin{align}
 \operatorname{Recall}_j@K&=\frac{\sum_{r=1}^{K_j}h_{jr}}{|G_j|},\\
 \operatorname{NDCG}_j@K&=
 \frac{\sum_{r=1}^{K_j}h_{jr}/\log_2(r+1)}
      {\sum_{r=1}^{\min(K,|G_j|)}1/\log_2(r+1)}.
\end{align}
The objective $\mathcal J$ is a macro mean over validation queries. By default it is mean
NDCG@20. The accepted objective names also include \code{recallK}, \code{ndcgK},
\code{ndcgfull}, \code{mrr} and \code{nll}, where $K$ in the first two names is a positive
integer. Full NDCG uses the ranks of all distinct ground-truth items. MRR averages the
reciprocal of the \emph{first} relevant item's rank.

For \code{nll}, first re-standardize each candidate score column using NumPy's population
standard deviation and epsilon $10^{-9}$; call it $\widetilde A_{ij}$. Define
$\ell_{ij}=\widetilde A_{ij}$ for non-seeds and $\ell_{ij}=-10^9$ for original seeds.
The higher-is-better objective is the negative of
\begin{equation}
 \operatorname{NLL}(A^{\mathrm{score}})=\frac1v\sum_j
 \left[\log\sum_{i\in\I}\exp(\ell_{ij})
       -\frac1{|G_j|}\sum_{i\in G_j}\ell_{ij}\right].
\end{equation}
The code uses a stable log-sum-exp calculation and float64 after this standardization.
The seed mask is the finite value $-10^9$, rather than literal negative infinity.

\subsection{Define observed evidence, prior and the final score}
Define an alias map $a(t)$ that replaces \code{simp} by \code{pair}, replaces
\code{diff2} by \code{diff2_pair}, and leaves every other signal key unchanged. For a
weight dictionary $w$, define the pair-only and full-incidence standardized fusions
\begin{equation}
 Y^{(0)}_{ij}=\sum_t w_t z^{a(t)}_{ij},\qquad
 Y^{(1)}_{ij}=\sum_t w_t z^t_{ij}.
\end{equation}
Given higher-order settings $(\alpha_h,\gamma_R)$, observed evidence is
\begin{equation}\label{eq:Y}
 Y_{ij}=(1-\alpha_h)Y^{(0)}_{ij}+\alpha_hY^{(1)}_{ij}
                         +\gamma_R z^{\sig{rao}}_{ij}.
\end{equation}
The $\alpha_h=0$ path uses only aliased pair signals. Increasing $\alpha_h$ mixes in
\emph{separately standardized} full signals. It does not simply add $\alpha_h$ times
raw incidence surplus to a fixed normalized pair score. Concordance enters through its
own standardized column and a single selected multiplier.

Before defining the prior, define the fallback key weight
$\widetilde w_t=w_t$ if key $t$ remains in the pruned dictionary, and
$\widetilde w_t=0.5$ otherwise. For prior-key list $\mathcal P$, define
\begin{equation}\label{eq:prior}
 P_{ij}=\sum_{t\in\mathcal P}\widetilde w_t z^t_{ij}
        +\gamma_R z^{\sig{rao}}_{ij}
        +\mathbf1[w_{\sig{pop}}>0]w_{\sig{pop}}z^{\sig{pop}}_{ij}.
\end{equation}
A missing popularity weight is interpreted as zero. The code includes the same concordance
term in the observed score and the prior. Popularity is also explicitly added to the prior
when selected, even though it was not one of the content/taxonomy prior candidates.

For nonnegative shrinkage $\kappa$ and scalar prior scale $s_0$, define the prior mean
$\mu_{0,ij}=s_0P_{ij}$ and candidate-dependent prior fraction
$u_i=\kappa/(f_i+\kappa)$ when $\kappa>0$. The final score is
\begin{equation}\label{eq:posterior}
 \widehat\theta_{ij}=\begin{cases}
 Y_{ij},&\kappa\le0,\\
 (1-u_i)Y_{ij}+u_i\mu_{0,ij},&\kappa>0.
 \end{cases}
\end{equation}
Thus a training-cold item ($f_i=0$) receives the scaled prior when shrinkage is enabled;
a high-exposure item receives more observed evidence. If selection disables shrinkage,
even a cold item follows $Y$. Scores can be negative and are not probabilities.

\subsection{Why this is called an empirical-Bayes router}
For interpretation, define latent compatibility $\theta_{ij}$, prior variance
$\sigma_0^2$ and observation-noise coefficient $\sigma_\epsilon^2$. The algebra of
\eqref{eq:posterior} matches the Gaussian assumptions
\begin{equation}
 \theta_{ij}\mid\mu_{0,ij}\sim\mathcal N(\mu_{0,ij},\sigma_0^2),\qquad
 Y_{ij}\mid\theta_{ij},f_i\sim
       \mathcal N(\theta_{ij},\sigma_\epsilon^2/f_i),\quad f_i>0,
\end{equation}
with ratio $\kappa=\sigma_\epsilon^2/\sigma_0^2$. The implemented extension to $f_i=0$
is the prior-only limit in \eqref{eq:posterior}. This is an interpretation of the scoring
rule; the code neither tests Gaussianity nor returns a posterior uncertainty distribution.
Moreover, the final ratio is selected from a validation grid, not forced to equal its
moment estimate.

\subsection{Stage A: coordinate search and backward pruning}
Initialize every menu weight to zero and the incumbent objective to $-10^{18}$. For at
most \code{fit_rounds} passes, visit the ordered menu. At coordinate $t$, try
$w_t\in\{0,0.25,0.5,1,1.5\}$ while holding other coordinates fixed. Skip an all-zero
candidate. Evaluate $\mathcal J(Y^{(0)})$, with concordance, incidence mixing and shrinkage
disabled. Accept a candidate only if its score exceeds the incumbent by $2\times10^{-4}$.
Stop after a pass with no accepted change. This is ordered greedy coordinate selection,
not an exhaustive search over all weight combinations.

Remove zero weights. Try removing each listed optional key, retaining at least one
positive-weight key; accept a removal if its objective is at least the current score minus
$3\times10^{-4}$. The visited pruning keys are:
\begin{center}\small
\begin{tabular}{@{}ll@{}}
\toprule Path & Optional keys, in order \\\midrule
General & \code{bridge}, \code{contrast}, \code{softcat}, \code{homo} \\
Research retail & \code{bridge}, \code{contrast}, \code{cat} \\
Research commerce & \code{contrast}, \code{homo} \\
\bottomrule
\end{tabular}
\end{center}
Absent keys, including unused \code{contrast} placeholders, are skipped. On the general
path, prior keys are chosen next; on baskets they are chosen after the higher-order grid.

\subsection{Stage B: choose prior keys and higher-order settings}
The ordered general research prior candidates are
\code{cat}, \code{softcat}, \code{style}, \code{imgstyle}, \code{bridge} for fashion and
\code{cat}, \code{homo}, \code{audiostyle}, \code{bridge} for music. The extended general path uses
\code{cat}, \code{softcat}, \code{homo}, \code{content}, \code{style}, \code{imgstyle}, \code{audiostyle}, \code{bridge}.
Keep available candidates with positive selected weights. If none remain, use all available
candidates in that order, allowing fallback weight $0.5$ in \eqref{eq:prior}.

Research retail uses positive-weight keys from \code{cat}, \code{content}, \code{bridge}, \code{contrast}; if
empty, its fallback is \code{cat}, \code{bridge}. Research commerce uses positive-weight keys
from \code{cat}, \code{homo}, \code{tfidf}, \code{bridge}, \code{pretrain}, \code{contrast}, \code{img}; its fallback is
\code{bridge}, \code{tfidf}, \code{cat}. Actual menus exclude the unimplemented placeholders.

Evaluate the score in \eqref{eq:Y} with $\kappa=0$, visiting $\alpha_h$ in
$\{0,0.5,1\}$ as the outer grid and $\gamma_R$ as the inner grid. The precise rules are:
\begin{center}\small
\begin{tabularx}{\linewidth}{@{}p{.19\linewidth}p{.18\linewidth}Y@{}}
\toprule Path & Concordance grid & Acceptance and incumbent \\\midrule
General & $0,0.2,0.35$ & Start with $(0,0)$ and its actual post-pruning score. Replace only for improvement $>5\times10^{-5}$. If concordance is unavailable, use only $\gamma_R=0$. \\
Retail & $0,0.15,0.25$ & Start objective at $-10^{18}$. Replace when candidate score is $\ge$ incumbent minus $10^{-5}$; a later near-tie can win. \\
Commerce & $0,0.2,0.35$ & Start objective at $-10^{18}$. Replace only for improvement $>10^{-5}$. \\
\bottomrule
\end{tabularx}
\end{center}
The retail rule permits small downward steps between successive accepted near-ties. It is
not equivalent to taking the first global argmax. The basket function subsequently carries
$\max(\text{coordinate-search incumbent},\text{accepted higher-order score})$ into its
shrinkage comparison. Unlike the general selector, it does not simply reset that baseline
to the final pruned higher-order score. These details are part of the current execution.

\subsection{Stage C: freeze prior scale and estimate exposure moments}\label{sec:moments}
With selected weights, prior keys and higher-order settings, flatten $Y$ and $P$ over
\emph{all} validation queries and catalogue items. Let $\operatorname{sd}_\nu$ be standard
deviation with correction $\nu$ on this flattened population. The scale is
\begin{equation}
 s_0=\begin{cases}
 \dfrac{\operatorname{sd}_1(Y)+10^{-9}}
       {\operatorname{sd}_1(P)+10^{-9}},&\text{general streaming path},\\[6pt]
 \dfrac{\operatorname{sd}_0(Y)}
       {\operatorname{sd}_0(P)+10^{-9}},&\text{NumPy research baskets}.
 \end{cases}
\end{equation}
This one scalar is fitted once and frozen. It is not re-estimated on a test batch.
The general selector computes global sums and squared sums in float64 across cached
blocks; standard deviations are not averaged block by block.

For each warm item define its validation mean squared residual $e_i$, inverse exposure
$x_i$, and their warm-item means $\overline e,\overline x$:
\begin{equation}
 e_i=\frac1v\sum_{j=1}^v(Y_{ij}-s_0P_{ij})^2,
 \qquad x_i=1/f_i,\qquad f_i>0.
\end{equation}
These are residuals between two score constructions, not squared errors against binary
relevance labels. The raw ordinary-least-squares estimates of slope $b_0$ and intercept
$a_0$ are
\begin{equation}
 b_0=\frac{\sum_i(x_i-\overline x)(e_i-\overline e)}
            {\sum_i(x_i-\overline x)^2},\qquad
 a_0=\overline e-b_0\overline x,
\end{equation}
where sums run over warm items. If fewer than two warm items exist or the denominator is
$\le10^{-12}$, the components are marked unidentifiable and $\widehat\kappa=0$.
Otherwise define the admissible plug-in quantities
\begin{equation}
 \widehat\sigma_0^2=\max(a_0,10^{-4}),\quad
 \widehat\sigma_\epsilon^2=\max(b_0,10^{-4}),\quad
 \widehat\kappa=\operatorname{clip}
 \left(\frac{\widehat\sigma_\epsilon^2}{\widehat\sigma_0^2},0.05,50\right).
\end{equation}
Raw estimates and truncation/identifiability flags are retained. These are clipped OLS
moment estimates, not an iterative constrained-optimization fit. The general streaming
path forms residuals and squares in the score dtype before float64 accumulation; NumPy
basket estimation converts its input score arrays to float64 before forming residuals.
One should not infer bit identity between those two numerical paths.

\subsection{Stage D: choose the shrinkage setting}
Use the following grids for $\kappa$, in the stated order:
\begin{center}\small
\begin{tabular}{@{}ll@{}}
\toprule Path & Candidate values \\\midrule
General & $0,0.05,0.5,1,\widehat\kappa,5$ \\
Research retail & $0,0.05,0.25,\widehat\kappa,1$ \\
Research commerce & $0,0.1,0.5,\widehat\kappa,2$ \\
\bottomrule
\end{tabular}
\end{center}
Unidentifiable variance components restrict every grid to the single value zero, even if
positive constants were otherwise listed. The general selector tries the fitted scale
$s_0$ and then scale $1$, using the ratio grid inside each scale. It accepts objective
improvements $>5\times10^{-5}$. Basket selectors keep the fitted scale and accept
improvements $>10^{-4}$ relative to their carried incumbent. If scale 1 wins on the general path, it becomes the frozen $s_0$; the recorded moment
estimates remain those fitted using the originally estimated scale. All start with shrinkage
disabled. A grid may contain a repeated value; the code does not require deduplication.

The resulting frozen state is the ordered menu, pruned weight dictionary, prior-key list,
$(\alpha_h,\gamma_R)$, $(\kappa,s_0)$, variance diagnostics, graph/content artifacts,
and dataset provenance. Test evaluation never changes these selected values.

\subsection{How validation memory and replay work}
The general path calls the independent all-signal scorer in blocks of
\code{validation_chunk}, default 32. It writes each required standardized signal to a
query-major float32 memory-mapped array with shape $(v,n)$. The approximate cache payload
for $d_z$ retained signal keys is $4vnd_z$ bytes, excluding headers. All validation queries
remain in the selection population; the block size limits resident scoring work, not the
number of validation examples. Candidate scores are macro-averaged with query-count
weights across blocks. Objective values are memoized by weights, higher-order settings,
backoff and prior-key order. Independent candidates can share one pass through the cache.

The cache is temporary and is closed/deleted after selection, including on exceptions.
It is not a fitted checkpoint. Research baskets instead compute their complete validation
signal matrices in memory and call the NumPy selector; \code{validation_chunk} does not
bound this basket validation allocation in the current code.

\section{Step 9: save fitted state and restore it exactly}\label{sec:checkpoint}
\src{bundlefuse/model.py: save, load_model; evaluator.py; validate.py}
\code{fit} marks the object fitted only after selection succeeds. It does not save a file.
A later \code{save(path)} makes a shallow model snapshot and a separate graph snapshot.
For baskets it also copies the provider and rebinds that provider to the copied graph.
The live fitted graph is not modified by checkpoint storage.

With \code{compressed=true}, the copied graph omits the derived attributes
\code{C}, \code{C_bin}, \code{PPMI}, \code{W_pair_sq}, \code{W_csc}, and \code{P_rw}.
It retains incidence matrices, exposure/size vectors, $W$, spectral features, selected
weights/routing, content/context artifacts and provenance. The file is Python pickle
with highest available protocol, despite its \code{.pt} extension. It is not a gzip archive,
and compression here does not mean quantizing weights or changing retained numeric values.
The optional fit-time content projection in Section~\ref{sec:features} is a separate,
score-changing approximation; current dataset JSONs disable that projection.

The snapshot is written to a temporary file in the destination directory, then atomically
replaces the destination. Temporary partial-write files are cleaned up; the final checkpoint
is retained. The ordinary evaluator's default destination is
\code{checkpoints/<dataset>.pt}, so a later save to that same name replaces it.
\code{validate.py} instead allocates a unique dataset/protocol/seed/UUID directory under
\code{checkpoints/validation/}, or under the supplied checkpoint root.

On inference, \code{ensure_graph_matrices} reconstructs missing $C$ from $B^\top B$
with its diagonal removed, and reconstructs $M$ using exactly \eqref{eq:itempmi} when
requested. These reconstructed matrices are retained on that loaded graph. It does not
rebuild every omitted helper artifact. In particular, concordance converts retained $W$
to CSC format if \code{W_csc} is missing, but keeps the conversion in a local variable:
the current routine can repeat that conversion on later calls. Core prediction uses the
factored $L$, so it does not need stored \code{P_rw} or \code{W_pair_sq}.

The custom unpickler remaps the earlier basket-provider module name to
\code{bundlefuse.basket_signals.BasketSignals}. Loading an ordinary checkpoint does not
refit SVD, clustering, weights, prior scale or shrinkage.

The fit CLI checks a validation-score probe before save, after save on the live model,
and after reload; the comparison is NumPy array-value equality. The validation runner
additionally compares every requested test score through its float32 unsigned-integer
view, giving an actual bit-pattern equality check. Both write checkpoint paths/hashes and
sizes to run reports. The temporary signal cache and atomic-write staging files are the
objects cleaned up automatically, not the validation runner's completed model files.

\section{Step 10: execute a new prediction}\label{sec:predict}
\src{bundlefuse/model.py: score_batch, score, predict; bundlefuse/selected_signals.py; bundlefuse/metrics.py}

\subsection{Validate the query and choose scoring seeds}
A fitted model accepts nonempty queries of distinct integer IDs in $[0,n)$; an explicit
user list must align with the query count. Define the seed-cap setting $c_{\max}\ge0$.
The prepared scoring context is
\begin{equation}
 S^\star=\begin{cases}
 S\text{ in its supplied order},&c_{\max}=0,\\
 \text{first }\min(c_{\max},|S|)\text{ IDs in sorted }S,&c_{\max}>0.
 \end{cases}
\end{equation}
This preparation is also used when computing validation signals. Original $S$ remains
available to metric/ranking code. Thus capped-away revealed items are neither scoring
seeds nor recommendation candidates. Seed capping does not turn them into ground truth;
the separate stress experiment below does change ground truth.

\subsection{Determine dependencies of the frozen fusion}
Let $\mathcal W=\{t:w_t\ne0\}$ be the remaining positive-weight keys. Before computing
scores, define the required set
\begin{align}
 \mathcal D={}&\{a(t):t\in\mathcal W\}\\
 &{}\cup\begin{cases}\mathcal W,&\alpha_h\ne0,\\\varnothing,&\alpha_h=0\end{cases}
 \cup\begin{cases}\{\sig{rao}\},&\gamma_R\ne0,\\\varnothing,&\gamma_R=0\end{cases}\nonumber\\
 &{}\cup\begin{cases}
 \{t\in\mathcal P:\widetilde w_t\ne0\}
 \cup\{\sig{pop}:w_{\sig{pop}}>0\},&\kappa>0,\\
 \varnothing,&\kappa\le0.
 \end{cases}\nonumber
\end{align}
This is \code{required_fusion_signals}. General \code{SelectedSignals} and basket
\code{score_selected} compute exactly those keys and their operator dependencies.
Indicators and overlap products are shared within a call; zero-weight/unneeded score
families are skipped. Candidate items are not pruned: this still scores the full catalogue.

For example, if only \code{simp} has positive weight, $\alpha_h=0$, $\gamma_R=0$ and
$\kappa=0$, the required signal is only \code{pair}; incidence surplus, concordance,
content and taxonomy are unnecessary. With $\alpha_h>0$, both \code{pair} and
\code{simp} are required. With $\kappa>0$, any fallback-prior keys are required too.

\subsection{Compute, standardize, fuse and rank}
\begin{enumerate}
\item Restore only missing graph matrices required by $\mathcal D$.
\item Build indicators for $S^\star$, compute required raw scores from
      Section~\ref{sec:signals}, and standardize them with \eqref{eq:z}.
\item Apply \eqref{eq:Y}, then \eqref{eq:prior} and \eqref{eq:posterior} when shrinkage is enabled.
      General scoring returns an $n\times q$ float32 NumPy array; basket fusion also returns NumPy scores.
\item For top-$K$ prediction, copy each score column and assign $-\infty$ to \emph{all}
      original revealed IDs $S$. Clamp output length to the number of available non-seed items.
\item Use partitioning to find the top-$K$ boundary. Include all strictly higher items,
      fill boundary ties in ascending ID order, then sort the selected candidates by
      descending score and ascending ID.
\end{enumerate}
\code{score_batch} returns unmasked full-catalogue scores; \code{score} extracts one
column; \code{predict} adds seed exclusion and deterministic top-$K$ selection.
\code{reference=True} uses the all-signal reference scorer before the same fusion.
The code's optimized and reference paths are intended to preserve operation order,
while general floating-point batching can still select different low-level matrix kernels.
Per-query standardization and frozen prior scale introduce no dependence on \emph{other}
query columns, but this is not a claim of bit equality across all BLAS kernels or batch shapes.

\section{Step 11: run evaluation and bind it to the fitted data}\label{sec:evaluate}
\src{inference.py; bundlefuse/cli_data.py; bundlefuse/evaluation.py; bundlefuse/provenance.py}
\code{inference.py} has three modes: direct/query-from-test prediction, test evaluation
at a requested Recall/NDCG cutoff, and a fuller JSON report. No mode refits the model.
The loaded object's model type and recorded dataset are checked against the requested CLI
arguments. Direct \code{--seeds} prediction can run from a checkpoint alone; prediction
with \code{--from-test}, evaluation and reporting reconstruct a dataset.

\subsection{Recorded fingerprints and actual acceptance checks}
The dataset specification records dataset name, item count, profile, split seed/mask,
query limits and totals, complete/partial-fold flags, original query IDs, and hashes of
training rows, validation/test pairs and available input files. Row hashes are SHA-256 of
compact JSON serialization in the actual row order.

For checkpoint-bound reconstruction, omitted CLI options inherit the checkpoint's protocol,
profile, split seed, mask and validation limit. A different dataset/profile/protocol, or a
changed resplit seed/mask, is rejected. A legacy checkpoint missing the profile field is
interpreted as extended. The loader then compares \code{n_items},
\code{train_sha256}, \code{val_pairs_sha256} and \code{input_sha256} with fitted
metadata and rejects mismatches. An empty validation or test set is also rejected.

The test pair hash is recorded but is not in that four-key equality loop. Test query caps
can be changed, for example to inspect a prefix through \code{--from-test}. Direct
explicit-seed prediction does not read and re-hash the raw dataset files. This describes
the implemented checks, rather than assuming that every stored field is always compared.
Runtime metadata additionally hashes Python package sources and the three entry points,
and records Python/platform/library versions and Torch thread count. Such hashes identify
a run; ordinary inference does not reject every change in runtime source hash.

\subsection{Macro test metrics and seed experiments}
Metrics use the ranking and equations in Section~\ref{sec:objectives}. Evaluation scores
query blocks, computes each block's macro metrics, multiplies by block query count, and
divides the total by the complete number of evaluated queries. A small final block
therefore has its correct weight. The standard full report includes cutoffs 10, 20 and 50.

Two different analyses appear in \code{report}:
\begin{description}[style=nextline]
\item[Observed seed-count cohorts]
Group existing test pairs by original $|S_j|$ and evaluate those unchanged pairs. The number
of queries per cohort is reported; the model's seed cap can further restrict actual context.
\item[Controlled seed stress]
For requested count $c\in\{1,2,3\}$, let $S_j^{(c)}$ be the first $\min(c,|S_j|)$ sorted
seeds and define
\begin{equation}
 G_j^{(c)}=G_j\cup(S_j\setminus S_j^{(c)}).
\end{equation}
Evaluate the new pair $(S_j^{(c)},G_j^{(c)})$. Removed seeds become targets, preserving total
bundle membership; the number of targets can increase. No fitting is repeated. This is a
new completion task, not merely a timing shortcut for the original test pair.
\end{description}
\code{validate.py} is a separate integration runner: it fits, saves permanently, reloads,
checks every requested test-score bit pattern, and records per-query metrics. It does
not prove equality with an external published result unless an independent matched
reference comparison is performed separately.

\section{A small numerical walk-through}\label{sec:example}
This example illustrates the equations; its frozen fusion settings are supplied for
explanation, not claimed to be settings selected by the validation routine. Let $n=5$,
training bundles be $b_1=\{0,1,2\}$ and $b_2=\{0,3\}$, and let original seeds be
$S=S^\star=\{0,1\}$. Item 4 is in the known catalogue but never appears in training.
Use the default exponents and $\beta=1.5$. Then
\begin{equation}
 B=\begin{pmatrix}1&1&1&0&0\\1&0&0&1&0\end{pmatrix},\qquad
 \bm f=(2,1,1,1,0)^\top.
\end{equation}
Define $a=1/\sqrt7$ and $b=1/\sqrt6$. The weighted overlaps and counts are
\begin{equation}
 \bm o=\left(\frac{a+b}{2},\frac{a}{\sqrt3}\right)^\top,
 \qquad \bm k=(2,1)^\top.
\end{equation}
The common pair score and surplus are therefore
\begin{align}
 \bm p&=(o_1+o_2,o_1,o_1,o_2,0)^\top,\\
 \bm\delta&=(1.5o_1,1.5o_1,1.5o_1,0,0)^\top.
\end{align}
Only the first training bundle receives incidence activation, because only it overlaps
both seeds. Suppose the frozen settings are $w_{\sig{simp}}=1$, $\alpha_h=1$,
$\gamma_R=0$, prior keys $\{\sig{content}\}$, $\kappa=1$, and $s_0=1$.
The prior key lacks a selected weight, so \eqref{eq:prior} assigns it weight $0.5$.
For illustration, supply raw content scores proportional to $(0,1,2,3,4)^\top$; their
Torch-standardized column is approximately
\begin{equation}
 \bm z^{\sig{content}}=(-1.265,-0.632,0,0.632,1.265)^\top.
\end{equation}
The observations are $Y=\bm z^{\sig{simp}}$, obtained by standardizing
$\bm p+\bm\delta$ across all \emph{five} catalogue items, including both seeds and item 4.
The warm candidate weights are $u_2=u_3=1/2$, while $u_4=1$. Consequently
\begin{equation}
 \widehat\theta_2=\tfrac12 z^{\sig{simp}}_2,\qquad
 \widehat\theta_3=\tfrac12 z^{\sig{simp}}_3+\tfrac14 z^{\sig{content}}_3,\qquad
 \widehat\theta_4=\tfrac12 z^{\sig{content}}_4\simeq0.632.
\end{equation}
The resulting candidate scores are approximately $\widehat\theta_2=0.288$,
$\widehat\theta_3=-0.273$, and $\widehat\theta_4=0.632$. For $K=2$, the returned
IDs are therefore $[4,2]$. Original IDs 0 and 1 are excluded only at ranking. This shows how exposure controls the
mixture: item 4 can receive content-prior evidence despite zero collaborative exposure,
while item 2 receives the bundle-overlap surplus. The model does not infer a new graph
edge or retrain a network to produce this cold-item score.

\appendix
\section{Extended-profile behavior and additional signals}\label{app:extended}
The extended profile uses the general model/scorer with a larger signal pool and different
basket preprocessing. It is part of the current code and helps explain checkpoints fitted
under earlier configurations; it is not interchangeable with the default research basket path.

\subsection{Preprocessing differences}
Extended retail reads the same capped/filterable basket pool but deduplicates by
\emph{household plus item set}, rather than item set alone. Missing product metadata share
a common fallback category. TF--IDF/SVD uses the shared 8,000-feature helper; the model
later normalizes its dense output. It builds training household histories and query users,
and provides a weighted auxiliary household incidence to the general model.

Extended commerce uses joined title text only with the 8,000-feature helper, treats the
category CSV field as a literal label rather than parsing category paths, and reads external
pretraining interactions as user histories. It does not construct the research commerce
session graph or two-hop sparse text bridge. Shared split/masking rules still apply.
The profile settings remain independent: crossing a research loader with an extended scorer,
or the reverse, follows the dispatch condition in Section~\ref{sec:dispatch}; it does not
recreate the paired research recipe automatically.

\subsection{Auxiliary co-occurrence and group signals}
Let $Z\in\R^{n_u\times n}$ be a supplied context-by-item incidence. The actual scorer
uses $Z^\top Z\bm x$, or $Z^\top ZD\bm x$ when its weighted-seed flag is set. It does
not remove the resulting diagonal. The basket loader's shared helper first normalizes
context rows by $(\sum_i Z_{ui}+1)^{-1/2}$. Extended retail uses this for household context
with weighted seeds; extended commerce uses it for pretraining with unweighted seeds.

The optional \code{group_labels} API instead constructs item-by-group membership $H_g$,
normalizes each group column by the square root of its catalogue size plus one, and sums
$\omega_g\widetilde H_g\widetilde H_g^\top\bm x$ over supplied groupings with their
provided weights $\omega_g$. This supplies non-music \code{homo}. Without such groups,
the general non-music fallback is raw category-count co-occurrence $A^\top A$, diagonal
kept, scored using the seed labels. This fallback is not the commerce category-PMI formula.

\subsection{Optional personalization}
Let $R_{iu}$ be the number of times item $i$ occurs in the supplied history of user $u$,
and let $\bm h_u$ be its binary item indicator. The general path supplies
\begin{equation}
 r^{\sig{userfreq}}_i(S,u)=R_{iu},\qquad
 \bm r^{\sig{usercooc}}(S,u)=C\bm h_u.
\end{equation}
These have code keys \code{user_freq} and \code{user_cooc}. An unknown/omitted user gives
zero columns. The scoring seed set does not define the history. Extended retail histories
use training baskets; extended commerce histories come from its supplied external
pretraining interactions. The research basket provider ignores per-query user arguments
and instead uses its shared training context graphs.

\subsection{Extended menu order}
For the general path, computed keys begin with \code{simp, diff2, pair, diff2_pair, pop},
then available \code{rao}, \code{cooc}, \code{ppmi}, and auxiliary keys in supplied dictionary
order. They continue with \code{cat}, available \code{homo}, available \code{softcat},
feature keys in construction order, then available user keys. Feature construction begins
with raw \code{content}, then styles in modality order, then \code{bridge}, then
\code{csvd}; with no modalities it contains only \code{csvd}. The selectable menu filters
out auxiliary mixture keys and disallowed families. Thus feature order and profile are
relevant inputs to the ordered coordinate search, even when two paths contain similar
mathematical score families.

\section{Baseline paths implemented beside BundleFuse}\label{app:baselines}
\src{bundlefuse/baselines.py}
These baselines share the CLI and ranking/metric functions, but do not run BundleFuse's
fusion search. Their graph fit constructs a common graph and retains incidence and
frequency vectors. Define $A_B=B^\top B$ \emph{with diagonal retained}, and define
$g_i$ as the zero-diagonal co-occurrence row mass plus $10^{-9}$ from
\eqref{eq:itempmi}. Define $D_1=\diag((f_i+1)^{-1/2})$ and
$G_1=\diag(g_i^{-1})$. The implemented scores are
\begin{center}\small
\begin{tabularx}{\linewidth}{@{}p{.23\linewidth}Y@{}}
\toprule Name & Actual score computation \\\midrule
\code{popularity} & $\log(1+\bm f)$, independent of seeds. \\
\code{cooccurrence} & $A_B\bm x$. \\
\code{itemknn} & $D_1A_BD_1\bm x$, with no explicit neighbor truncation. \\
\code{p3alpha} & Three iterations $\bm w\leftarrow G_1A_B\bm w$, starting at $\bm x$. \\
\code{rp3beta} & The preceding three-step output divided by $(f_i+1)^{0.5}$. \\
\code{contentknn} & Dot product of normalized concatenated content with the mean normalized seed feature; no explicit neighbor truncation. \\
\bottomrule
\end{tabularx}
\end{center}
The names are interface labels: the implemented \code{p3alpha} numerator retains a
self-diagonal while its denominator uses off-diagonal degree, so it is not literally a
row-stochastic walk on the matrix used in the numerator. No tunable alpha exponent is
applied by that class. These exact baselines should not be silently identified with every
published variant bearing a similar name. Their supplied features and data folds still
follow the same chosen loader profile.

\section{Runtime cost and the timing entry point}\label{app:runtime}
\src{scripts/benchmark_inference.py; bundlefuse/operators.py: build_indicator, score_rao_concordance}
Let $d$ be the dense feature width, $k=|S^\star|$, and $\nnz(A)$ the stored nonzero
count of sparse matrix $A$. Dense candidate dot products require approximately $O(nd)$
work per query after forming the seed mean. Sparse incidence/co-occurrence operations
have costs driven by their nonzero structure. A dense score block requires $O(nq)$
memory per simultaneously materialized score family. Retaining item vectors does not
turn this implementation into an approximate nearest-neighbor index: candidates are
still scored over all $n$ items.

Concordance forms an $n\times k$ support matrix; its dense-feature multiplication costs
$O(ndk)$ and its direct support--kernel product costs $O(nk^2)$. A missing CSC copy of
$W$ introduces conversion/storage work on the order of $\nnz(W)+n$ per concordance call
in the current implementation. The query indicator currently starts as an $n$-row LIL
matrix before conversion to CSR, which also has allocation overhead. These are actual
execution costs, not extra mathematical terms in the score.

The benchmark script hashes and loads retained checkpoints, reconstructs/fingerprints
queries, performs one first request, warms up, and times the public \code{predict(k=20)}
call on sampled distinct test queries in shuffled repetitions. Its default sample is 256
queries, three repetitions, up to 16 warm-up queries, and one thread for Torch and native
thread pools. Dataset workers run sequentially in fresh processes. Timed online work
includes input checks, scoring, fusion, exclusion and top-$K$ ranking; it excludes
checkpoint/query loading, warm-up, metric evaluation and consistency hashing. A separate
batch-size-32 timing includes both \code{score_batch} and per-column top-$K$. It reports
amortized batch throughput separately from single-request latency, and does not modify
checkpoints. This explains what a timing measurement of this code represents without
importing empirical performance numbers into the theory.

\section{Active pipeline versus helper methods}\label{app:inactive}
The call path above uses incidence surplus, \code{score_rao_concordance}, dense/sparse
content bridges, validation selection and conjugate score mixing. Several other helpers
exist in \code{operators.py} but are not called by the described \code{BundleFuse.fit}
and ordinary prediction path: \code{propagate_rw}, \code{score_concordance_e2},
\code{score_simplicial_and_diff2}, \code{build_cold_warm_knn},
\code{fit_eb_moments}, \code{apply_tail_cold_calibration}, and
\code{precision_backoff_weights}. In particular, this pipeline does not insert cold items
into reserved top-$K$ slots, perform an extra KNN cold-item transfer, or add the separate
unweighted elementary-symmetric concordance helper. Their presence in the file is not
evidence that those mechanisms execute. The theoretical interpretation must follow the
caller, not every function name found in the module.

\section{Source map and reproduction}\label{app:map}
Each row below identifies the code responsible for a stage. Paths are relative to the
repository root. The companion JSON manifest supplies full file hashes; no previous
LaTeX notes or paper sections are inputs to the derivation.
\small
\begin{longtable}{@{}>{\raggedright\arraybackslash}p{.26\linewidth}>{\raggedright\arraybackslash}p{.68\linewidth}@{}}
\toprule Stage & Source / callable \\\midrule\endhead
Configuration and registry & \code{bundlefuse/config.py}: \code{make_config}; \code{bundlefuse/datasets/__init__.py}: \code{DatasetSpec}. \\
Provided folds and masks & \code{bundlefuse/splits.py}: \code{ProvidedSplits}, \code{mask}. \\
Data objects and full-ranking load & \code{bundlefuse/datasets/base.py}: \code{Dataset.fit_kwargs}, \code{load_bundlemllm}. \\
Research basket load & \code{bundlefuse/datasets/baskets.py}: \code{load_retail}, \code{load_commerce}. \\
Graph and collaborative features & \code{bundlefuse/operators.py}: \code{SimplicialGraphOperator.__init__}, \code{propagate}, \code{score_pairwise_and_diff2}. \\
Higher-order scores & Same class: \code{score_simplicial_incidence_only}, \code{score_rao_concordance}. \\
Content and ontology & \code{bundlefuse/operators.py}: \code{BilateralContentBridge}, \code{DualSpectrumOntology}; \code{bundlefuse/model.py}: \code{BundleFuse.fit}. \\
General full/selected scores & \code{bundlefuse/model.py}: \code{_score_all}; \code{bundlefuse/selected_signals.py}: \code{SelectedSignals.score}. \\
Research basket scores & \code{bundlefuse/basket_signals.py}: \code{BasketSignals}, \code{score_selected}. \\
Standardization & \code{bundlefuse/operators.py}: \code{zcols}. \\
Fusion and dependency selection & \code{bundlefuse/fusion.py}: \code{apply_fusion}, \code{apply_basket_fusion}, \code{required_fusion_signals}. \\
Validation selection & Same module: \code{fit_fusion}, \code{fit_basket_fusion}, \code{ValidationSignals}. \\
Validation storage & \code{bundlefuse/validation_cache.py}: \code{ValidationCache}. \\
Exposure moments and posterior & \code{bundlefuse/operators.py}: \code{EmpiricalBayesRouter}. \\
Public inference / checkpoint & \code{bundlefuse/model.py}: \code{score_batch}, \code{score}, \code{predict}, \code{save}, \code{load_model}. \\
Ranking and metrics & \code{bundlefuse/metrics.py}: \code{top_items}, \code{ranking_metrics}, \code{full_ndcg}, \code{mean_reciprocal_rank}. \\
Dataset binding and reports & \code{bundlefuse/cli_data.py}: \code{load_data}; \code{bundlefuse/evaluation.py}; \code{bundlefuse/provenance.py}. \\
Executables & \code{evaluator.py}, \code{inference.py}, \code{validate.py}, \code{scripts/benchmark_inference.py}. \\
\bottomrule
\end{longtable}
\normalsize
To compile this standalone source from the repository root:
\begin{lstlisting}
pdflatex -interaction=nonstopmode -halt-on-error \
  -output-directory=output/pdf docs/pipeline_theory.tex
pdflatex -interaction=nonstopmode -halt-on-error \
  -output-directory=output/pdf docs/pipeline_theory.tex
\end{lstlisting}
Run twice to resolve the table of contents and cross-references. The generated PDF is
\code{output/pdf/pipeline_theory.pdf}. The document and companion manifest describe the
inspected working-tree snapshot; later code changes require checking the equations and
updating the source hashes rather than merely changing a document date.
\end{document}
